\documentclass{midl} % Include author names

% The following packages will be automatically loaded:
% jmlr, amsmath, amssymb, natbib, graphicx, url, algorithm2e
% ifoddpage, relsize and probably more
% make sure they are installed with your latex distribution

\usepackage{mwe} % to get dummy images

% Header for short paper
\jmlrproceedings{}{Methods in AI research}
\jmlrpages{}
\jmlryear{2026}
\jmlrworkshop{Short Paper}

\title[Short Title]{Full Title of Article}

 % Use \Name{Author Name} to specify the name.
 % If the surname contains spaces, enclose the surname
 % in braces, e.g. \Name{John {Smith Jones}} similarly
 % if the name has a "von" part, e.g \Name{Jane {de Winter}}.
 % If the first letter in the forenames is a diacritic
 % enclose the diacritic in braces, e.g. \Name{{\'E}louise Smith}

\midlauthor{
\Name{Author Name1} \Email{abc@sample.edu}\\
\Name{Author Name2} \Email{xyz@sample.edu}\\
\Name{Author Name3} \Email{alphabeta@example.edu}\\
\Name{Author Name4} \Email{uvw@foo.ac.uk}\\
}

\begin{document}

\maketitle

\begin{abstract}
This is a great paper and it has a concise abstract.
\end{abstract}

\section{Introduction}

This is where the content of your paper goes.  Some random notes:
\begin{itemize}
\item You should use \LaTeX \cite{Lamport:Book:1989}.
\item JMLR/PMLR uses natbib for references. For simplicity, here, \verb|\cite|  defaults to
  parenthetical citations, i.e. \verb|\citep|. You can of course also
  use \verb|\citet| for textual citations.
\item You should follow the guidelines provided by the conference.
\item Read through the JMLR template documentation for specific \LaTeX
  usage questions.
\item Note that the JMLR template provides many handy functionalities
such as \verb|\figureref| to refer to a figure,
e.g. \figureref{fig:example},  \verb|\tableref| to refer to a table,
e.g. \tableref{tab:example} and \verb|\equationref| to refer to an equation,
e.g. \equationref{eq:example}.
\end{itemize}

\begin{table}[htbp]
 % The first argument is the label.
 % The caption goes in the second argument, and the table contents
 % go in the third argument.
\floatconts
  {tab:example}%
  {\caption{An Example Table}}%
  {\begin{tabular}{ll}
  \bfseries Dataset & \bfseries Result\\
  Data1 & 0.12345\\
  Data2 & 0.67890\\
  Data3 & 0.54321\\
  Data4 & 0.09876
  \end{tabular}}
\end{table}

\begin{figure}[htbp]
 % Caption and label go in the first argument and the figure contents
 % go in the second argument
\floatconts
  {fig:example}
  {\caption{Example Image}}
  {\includegraphics[width=0.5\linewidth]{example-image}}
\end{figure}

\begin{algorithm2e}
\caption{Computing Net Activation}
\label{alg:net}
 % older versions of algorithm2e have \dontprintsemicolon instead
 % of the following:
 %\DontPrintSemicolon
 % older versions of algorithm2e have \linesnumbered instead of the
 % following:
 %\LinesNumbered
\KwIn{$x_1, \ldots, x_n, w_1, \ldots, w_n$}
\KwOut{$y$, the net activation}
$y\leftarrow 0$\;
\For{$i\leftarrow 1$ \KwTo $n$}{
  $y \leftarrow y + w_i*x_i$\;
}
\end{algorithm2e}

\section{System architecture}

\section{Experiments}

\subsection{Error analysis}
In this section, we analyze the errors made by our systems. 
The systems we consider are the following: the rule-based baseline, 
a support vector machine and a decision tree classifier on a bag of words representation of the utterances and on frozen pretrained embeddings.
During the error analysis, we look at the groupled split only, because we do not want data leakage to infuence the results.
We focus on whether there are specific dialog acts and utterances that are more difficult to classify.
In addition, we note why these dialog acts and utterances are difficult to classify.\newline

First we look at the specific dialog acts that are more difficult to classify.
We do this by looking at the recall for each dialog act and each system.
\autoref{tab:dialog_acts_recall} shows the recall for each dialog act and each system.
We look at the recall, because we want to determine which dialog acts are difficult to classify.
This entails determining what proportion of the utterances of a class is correctly classified as that class.\newline

When we look at \autoref{tab:dialog_acts_recall}, we see that there are only a few dialog acts that 
are relatively well classified by all systems. These dialog acts are 'inform', 'request' and 'thankyou'.
These dialog acts have a relatively high number of training examples. 
This might help the systems learn more representative patterns for the dialog acts.\newline

The dialog acts 'reqalts', 'confirm' and 'bye' show an overall good performance, with the exception of the 
decision tree on the frozen pretrained embeddings. The dialog act 'negate' has the same overall good performance, besides the decision tree on the frozen pretrained embeddings, but it also 
has a relatively low recall on for the rule-based baseline. We see that the recall of the
decision tree on the frozen pretrained embeddings is relatively low for all dialog acts compared to the other systems.
This suggests that the lower recall is more likely related to the combination of 
a decision tree with the frozen pretrained embeddings.\newline

We see that the 'ack' and 'restart' dialog acts are diffucult to classify for all systems except for the rule-based baseline.
We are not able to say that these dialog act are difficult to classify, because the models were only tested on 3 and 2 utterances and 
trained on 23 and 5 utterances.
The rule-based baseline correctly classifed the dialog acts, this may be related to the manually defined rules
that directly capture the patterns of the dialog acts.\newline

The dialog acts 'repeat', 'hello' and 'deny' show a variation in recall between the systems.
However, we are not able to say anything about whether the dialog acts are difficult to classify, because the systems
were only tested on a few utterances. 
The 'reqmore' dialog acts has no test examples, so no conclusions can be drawn from the recall. \newline

The dialog act 'affirm' has a clear difference in recall for the systems. In this case, there are 103 utterances in the test set 
and 1032 utterances in the training set, so this difference cannot only be due to the small sets. 
This suggests that this dialog act might be difficult to classify for the support vector machine 
on a bag of words representation of the utterances and the support vector machine and decision tree on frozen pretrained embeddings.\newline

The only dialog act that performs poorly with every model and has a large enough test and training set is 'null'.
This dialog act has 57 test utterances and 1518 training utterances. It has a recall of 0.65 or lower for every system.
When we look at the utterances that are in the 'null' class, we see that there are fewer words that directly indicate that the utterance should be classified as 'null'.
This might make it more difficult to classify this dialog act correctly.\newline

The results show that the difficulty of classifying a dialog act depends on the system that is used to predict the dialog act.
In addition, the small training and test sets make it difficult to draw reliable conclusions from the results.
The recall for the dialog act 'null' is low across all systems. The recall shows that this dialog act is the most difficult to classify correctly for all systems.


\begin{table}[htbp]
\centering
\floatconts
  {tab:dialog_acts_recall}%
  {\caption{
    The recall for the different systems on the dialog acts, on the test set. BoW denotes bag-of-words representation, 
    BERT denotes frozen pretrained DistilBERT embeddings. 
    SVM denotes support vector machine, and DT denotes decision tree.
    The \#test and \#train are the total number of utterances that the specific dialog act has in the test set and the training set.
  }}%
  {\begin{tabular}{lccccccc}
  \bfseries Dialog act & \bfseries Rule-based & \bfseries DT BoW & 
  \bfseries SVM BoW & \bfseries DT BERT & \bfseries SVM BERT &
  \bfseries \#test & \bfseries \#train \\ 
  \hline
  ack       & 1.000 & 0.000 & 0.000 & 0.000 & 0.000 & 3 & 23 \\
  affirm    & 0.981 & 0.981 & 0.194 & 0.107 & 0.165 & 103 & 1032 \\
  bye       & 0.818 & 0.727 & 0.818 & 0.000 & 0.818 & 11 & 97 \\
  confirm   & 0.895 & 0.842 & 0.842 & 0.211 & 0.737 & 19 & 133 \\
  deny      & 0.500 & 0.500 & 0.500 & 1.000 & 0.000 & 2 & 24 \\
  hello     & 1.000 & 0.889 & 0.889 & 0.111 & 0.556 & 9 & 75 \\
  inform    & 0.907 & 0.952 & 0.991 & 0.721 & 0.994 & 2337 & 6978 \\
  negate    & 0.087 & 0.826 & 1.000 & 0.435 & 0.913 & 23 & 399 \\
  null      & 0.632 & 0.404 & 0.491 & 0.246 & 0.649 & 57 & 1518 \\
  repeat    & 0.333 & 0.667 & 0.667 & 0.000 & 1.000 & 3 & 29 \\
  reqalts   & 0.770 & 0.900 & 0.910 & 0.440 & 0.930 & 100 & 1535 \\
  reqmore   & 0.000 & 0.000 & 0.000 & 0.000 & 0.000 & 0 & 5 \\
  request   & 0.925 & 0.929 & 0.955 & 0.882 & 0.959 & 507 & 5721 \\
  restart   & 1.000 & 0.000 & 0.000 & 0.000 & 0.000 & 2 & 5 \\
  thankyou  & 1.000 & 0.700 & 0.850 & 0.850 & 0.900 & 20 & 3228 \\
  \end{tabular}}
\end{table}


For the dialog act 'null', we will look at some of the misclassified utterances:
\begin{itemize}
  \item \textbf{Text: how about} This was classified as 'reqalts', but should have been classified as 'null'.
        Sentences that start with "how about" are typically 'reqalts', but in this case the sentence is 
        left unfinished and should therefore be classified as 'null'. This makes classifying this utterance difficult for the systems.
  \item \textbf{Text: what about} This was also classified as 'reqalts', but should have been classified as 'null'. This is a 
        difficult utterance for the same reason as the preceding example: it is the start of a typical 'reqalts', but it is left unfinished.
  \item \textbf{Text: do they have} This was classified as 'request'. Again, this is the beginning of a typical 'request', but the sentence is left unfinished.
  \item \textbf{Text: is there any} This was classified as 'reqalts'. Again, this is the beginning of a typical 'reqalts', but the sentence is left unfinished.
  \item \textbf{Text: type} This was classified as 'inform'. In this case, the word "type" is a common word for 'inform', but here it forms the complete sentence. Therefore, the sentence does not have meaning and should be classified as 'null'.
        This makes classifying this as 'null' difficult.
\end{itemize}

The examples of wronly classified utterances indicate that unfinished sentences make the classification for the dialog act 'null' more difficult.
The utterances contain words that are commonly used in a specific dialog act, but are not used in a complete sentence. This way, the systems 
may classify the utterances based on the commonly used words, but miss that the sentence is not completed.


\section{Discussion}

\section{Conclusion}

\bibliography{midl-samplebibliography}

\appendix

\section{Data}

\section{System details}


\section{Configurability}

\section{Additional experiments}

\section{Contributions}

\section{AI disclosure statement}
We, Name1 (student number), Name2 (student number), Name3 (student number), and Name4 (student number) hereby confirm that the submitted work for the course \textit{Methods in AI research} is our own work. We declare that we have only used AI tools as permitted for this work. \\

\noindent Indicate whether you used AI tools for your assignment:
\begin{itemize}
    \item Yes, we used AI tools to help with the submitted work.
    \item No, we did not use AI tools to help with the submitted work.
\end{itemize}

\noindent If you selected "Yes", please indicate the following.\\

\noindent Which AI tool(s) did you use and for what purpose?

\begin{table}[htbp]
\floatconts
  {tab:ai_use}%
  {\caption{Overview of AI Tools Used}}%
  {\begin{tabular}{ll}
  \bfseries Name of the AI tool & \bfseries Purpose of using the AI tool\\
  e.g., Claude or Grammarly & Grammar check \\
   & \\
   & \\
   & \\

  \end{tabular}}
\end{table}

\noindent We acknowledge our responsibility as students/learners to thoroughly check all output and content produced by AI tools, and we accept full responsibility for its accuracy and validity.\\

\noindent Any comments:\\\\\\\\\\\\


\noindent Names and signatures of all students: \\\\\\\\\\\\



\noindent Date: 

\end{document}
